%\documentclass[12pt,oneside]{amsart}
\documentclass{scrartcl}

%\documentclass{article}

\usepackage{amsthm,amsmath,amssymb}
\usepackage[numbers]{natbib}

%\usepackage[colorlinks]{hyperref}
\usepackage{hyperref}
\hypersetup{
    colorlinks,%
    citecolor=black,%
    filecolor=black,%
    linkcolor=black,%
    urlcolor=black
}
\usepackage{hypernat}
\usepackage[top=2.5cm, bottom=2.5cm, left=2.8cm,
right=2.8cm]{geometry} 

%\setlength{\parskip}{0pt}
%\setlength{\parsep}{0pt}
%\setlength{\headsep}{0pt}
%\setlength{\topskip}{0pt}
%\setlength{\topmargin}{0pt}
%\setlength{\topsep}{0pt}
%\setlength{\partopsep}{0pt}

\linespread{1}

%\usepackage{marvosym}

%\textwidth = 6.5 in \textheight = 9.2 in \oddsidemargin = 0.0 in
%\evensidemargin = 0.0 in \topmargin = -0.0 in \headheight = 0.0 in
%\headsep = 0.0in \parskip = 0.0in \parindent = 0.0in
%\def\baselinestretch{0.871}


\newtheorem{theorem}{Theorem}[section]
\newtheorem{proposition}[theorem]{Proposition}
\newtheorem{problem}[theorem]{Problem}
\newtheorem{fact}[theorem]{Fact}
\newtheorem{lemma}[theorem]{Lemma}
\newtheorem{defn}[theorem]{Definition}
\newtheorem{corollary}[theorem]{Corollary}
\newtheorem{remark}{Remark}[section]
\newtheorem{example}[theorem]{Example}

\newcommand{\E}{{\mathbb E}}
\newcommand{\se}{{\mathcal{E}}}
\newcommand{\R}{{\mathbb R}}
\renewcommand{\P}{{\mathbb P}}
\newcommand{\ac}{{\mathcal{A}}}
\newcommand{\A}{{\mathcal{A}}}
\newcommand{\B}{{\mathcal{B}}}
\newcommand{\C}{{\mathcal{C}}}
\newcommand{\M}{{\mathcal{M}}}
\newcommand{\Mf}{{\mathfrak{M}}}
\newcommand{\T}{{\mathcal{T}}}
\newcommand{\Z}{{\mathcal{Z}}}
\newcommand{\K}{{\mathcal{K}}}
\newcommand{\lc}{{\mathcal{L}}}
\newcommand{\Ps}{{\mathcal{P}}}
\newcommand{\G}{{\mathcal{G}}}
\newcommand{\pa}{{\mathring{p}}}
\newcommand{\F}{{\cal F}}
\newcommand{\samp}{{\mathcal{X}}}
\newcommand{\X}{{\mathcal{X}}}
\newcommand{\hi}{{\hat{\Phi}}}
\newcommand{\data}{{\text{data}}}

\newcounter{rcnt}[section]
\renewcommand{\thercnt}{(\roman{rcnt})}

\def\qt#1{\qquad\text{#1}}

\def\argmin{\mathop{\rm argmin}}
\def\argmax{\mathop{\rm argmax}}


\setlength{\parskip}{1.4 \medskipamount}
%\sloppy
%\linespread{1.3}

\begin{document}

\title{STAT 238 - Bayesian Statistics  \\
  Lecture Four} 
\subtitle{Spring 2026, UC Berkeley}

\author{Aditya Guntuboyina}


\date{28 Jan 2026}

\maketitle

In the last lecture, we discussed the following problem. 

\section{Example 3: Inference from measurements}

\begin{problem}\label{lik3}
  Suppose a scientist makes 6 numerical measurements 26.6, 38.5, 34.4,
  34, 31, 23.6 on an unknown real-valued physical quantity
  $\theta$. On the basis of these measurements, what can be inferred
  about $\theta$?  
\end{problem}

\subsection{Frequentist Solution}

Here is the standard frequentist solution to this problem. Use the model:
\begin{align}
  \label{norm_iid}
  X_1, \dots, X_n \overset{\text{i.i.d}}{\sim} N(\theta, \sigma^2). 
\end{align}
It then follows that $\bar{X}_n := (X_1 + \dots + X_n)/n \sim
N(\theta, \sigma^2/n)$ which implies:
\begin{align}\label{std_norm}
  \frac{\sqrt{n}(\bar{X}_n - \theta)}{\sigma} \sim N(0, 1). 
\end{align}
If $\sigma$ is known, this gives the confidence interval $\bar{X}_n
\pm \frac{\sigma}{\sqrt{n}} z_{\alpha/2}$ for $\theta$ (here
$z_{\alpha/2}$ satisfies $\P\{N(0, 1) > z_{\alpha/2}\} =
\alpha/2$). But this confidence interval cannot be computed as
$\sigma$ is unknown. It is natural to replace $\sigma$ by the natural
estimator: 
\begin{align*}
  \hat{\sigma} := \sqrt{\frac{1}{n-1} \sum_{i=1}^n (X_i -
  \bar{X}_n)^2}. 
\end{align*}
But then the normal distribution in \eqref{std_norm} needs to be
changed to the Student $t$ distribution with $n-1$ degrees of freedom:
\begin{align}\label{std_t}
  \frac{\sqrt{n}(\bar{X}_n - \theta)}{\hat{\sigma}} \sim t_{n-1}. 
\end{align}
This leads to the confidence interval:
\begin{align}\label{freq_ci}
\left[  \bar{X}_n - \frac{t_{n-1,
      \alpha/2}}{\sqrt{n}}  
  \sqrt{\frac{1}{n-1} \sum_{i=1}^n (X_i - \bar{X}_n)^2}, 
  \bar{X}_n + \frac{t_{n-1,
      \alpha/2}}{\sqrt{n}}  
  \sqrt{\frac{1}{n-1} \sum_{i=1}^n (X_i - \bar{X}_n)^2}  \right]. 
\end{align}
When $\alpha = 0.05$ and one plugs in the observed data $X_1 = 26.6,
X_2 = 38.5, X_3 = 34.4, X_4 = 34, X_5 = 31, X_6 = 23.6$ with $n = 6$
in the above interval, we obtain the interval $[25.598, 37.102]$. 

\subsection{Bayesian Solution}
The Bayesian solution also leads to the same interval but with a
different reasoning. We went over the 
calculations in the last lecture. Here are the main facts. We use the
likelihood:  
\begin{equation}\label{likeli}
\text{Likelihood} = \prod_{i=1}^n \frac{1}{\sqrt{2 \pi} \sigma} \exp
\left(-\frac{(x_i - 
      \theta)^2}{2 \sigma^2} \right). 
\end{equation}
where $n = 6$ and $x_1 = 26.6, x_2 = 38.5, x_3 = 34.4,
  x_4 = 34, x_5 = 31, x_6 = 23.6$ denote the observed data
  points.

The unknown parameters are $\theta$ and $\sigma$. The prior is given by:     
\begin{align}\label{prior_norm}
  \theta, \log \sigma \overset{\text{i.i.d}}{\sim} \text{uniform}(-C, C)
\end{align}
for a very large positive constant $C$. In terms of the densities,
\eqref{prior_norm} is the same as
\begin{align*}
\text{prior} = f_{\theta, \sigma}(\theta, \sigma)
\propto \frac{\mathbf{1}\{-C < \theta < C,\,-C < \log \sigma < C\}}{\sigma}.
\end{align*}
For this model, the posterior becomes: 
\begin{align*}
  \text{posterior} = f_{\theta, \sigma \mid \text{data}}(\theta,
                     \sigma)\propto \sigma^{-n-1}
                     \exp\left(-\frac{1}{2 
    \sigma^2}\sum_{i=1}^n(x_i - \theta)^2 \right)\mathbf{1}\{-C < \theta < C,\,-C < \log \sigma <
    C\}. 
\end{align*}
This is the joint posterior density of $\theta$ and $\sigma$. The
posterior of $\theta$ alone is obtained by integrating out $\sigma$: 
\begin{align*}
  f_{\theta \mid \text{data}}(\theta) \propto \mathbf{1}\{-C <
  \theta < C\} \int_{e^{-C}}^{e^C} \sigma^{-n-1} \exp\left(-\frac{1}{2
    \sigma^2}\sum_{i=1}^n(x_i - \theta)^2 \right) d\sigma
\end{align*}
Because $C$ is large, the limits of the integral can be taken to be
$0$ and $\infty$. The integral then can be calculated precisely to
obtain 
\begin{align*}
  f_{\theta \mid \text{data}}(\theta) &\propto \mathbf{1}\{-C <
  \theta < C\} \left(\frac{1}{S(\theta)} \right)^{n/2} 
\end{align*}
where $S(\theta)$ is the sum of squares term:
\begin{align*}
  S(\theta) = \sum_{i=1}^n (x_i - \theta)^2. 
\end{align*}
If $C$ is large, then the indicator can be dropped (because it will
essentially be always 1) so the posterior becomes:
\begin{align*}
  f_{\theta \mid \text{data}}(\theta) &\propto
                                        \left(\frac{1}{S(\theta)}
                                        \right)^{n/2} \propto
                                        \left(\frac{S(\hat{\theta})}{S(\theta)}
                                        \right)^{n/2}. 
\end{align*}
where $\hat{\theta} = \bar{x} =
(x_1 + \dots + x_n)/n$ is the least squares estimator (which minimizes
$S(\theta)$ over all $\theta$). Thus the posterior mode is the mean.

It can be shown that this distribution is closely related to the
$t$-distribution. Specifically, 
\begin{align}\label{bayes_prob}
\frac{\sqrt{n}(\theta - \hat{\theta})}{\sqrt{S(\hat{\theta})/(n-1)}}
  \mid \text{data} \sim t_{n-1}
\end{align}
where $t_{n-1}$ denotes the $t$-density with $n-1$ degrees of
freedom. Note that $\hat{\theta} = \bar{x}$ and $S(\hat{\theta}) =
\sum_{i=1}^n (x_i - \bar{x})^2$.

So the Bayesian point estimate is simply $\hat{\theta} = \bar{x}$
(this is the posterior mean, median and mode!). A $100(1 - \alpha)$\%
uncertainty interval for $\theta$ is given by:
\begin{align}\label{cred_int}
\left[  \hat{\theta} - \frac{1}{\sqrt{n}} t_{n-1, \alpha/2}
  \sqrt{\frac{S(\hat{\theta})}{n-1}},   \hat{\theta} +
  \frac{1}{\sqrt{n}} t_{n-1, \alpha/2} 
  \sqrt{\frac{S(\hat{\theta})}{n-1}} \right] 
\end{align}
where $t_{n-1, \alpha/2}$ is the $(1 - \alpha/2)$ quantile of the
Student $t$-distribution with $n-1$ degrees of freedom. This
uncertainty interval is referred to as the Bayesian Credible Interval.   

Thus, in problem \ref{lik3}, the standard frequentist and Bayesian
solutions coincide.

\subsection{Frequentist vs Bayes}

However, it is very easy to break this coincidence. For example,
consider the following problem:
\begin{problem}\label{lik3_prime}
  A scientist gets repeated measurements on an unknown physical
  quantity $\theta$ until a measurement smaller than 25 is
  obtained. Suppose the resulting data is 26.6, 38.5, 34.4, 34, 31,
  23.6. On the basis of these measurements, what can be inferred about
  $\theta$?  
\end{problem}
Note that the observed data is exactly the same as before.

The
frequentist confidence interval \eqref{freq_ci} is no longer
valid, because the
frequentist probability statement \eqref{std_t} is no longer
valid. This is because the number of datapoints $n$ cannot be taken to
be deterministically equal to 6. So the frequentist probability that we
need to calculate is (below we denote the number of data points by $N$
and treat it as a random variable):
\begin{equation*}
  \P \left\{\bar{X}_N - \frac{t_{N-1,
      \alpha/2}}{\sqrt{n}}  
  \sqrt{\frac{1}{N-1} \sum_{i=1}^N (X_i - \bar{X}_N)^2} \leq \theta \leq
  \bar{X}_N + \frac{t_{N-1,
      \alpha/2}}{\sqrt{N}}  
  \sqrt{\frac{1}{N-1} \sum_{i=1}^N (X_i - \bar{X}_N)^2}  
\right\} 
\end{equation*}
where $X_1, X_2, \dots$ are i.i.d $N(\theta, \sigma^2)$ as before and
\begin{equation*}
  N := \inf \left\{n \geq 1 : X_n \leq 25 \right\}. 
\end{equation*}
The probability above is complicated and there is no reason for
it to be exactly equal to $1 - \alpha$. Constructing valid frequentist 
confidence intervals in the presence of \textit{stopping rules} (such
as the rule of stopping as soon as we observe a data point smaller than 25) 
is, in 
fact, a problem of current research (see e.g., the paper
\url{https://arxiv.org/abs/2210.01948}).

In contrast to frequentist inference, the Bayesian inference procedure
will not change. This is because the likelihood function in Problem
\ref{lik3_prime} is the same as the likelihood function in Problem
\ref{lik3}. To verify this, consider the following likelihood in Problem
\ref{lik3_prime} (below $\delta$ denotes the rounding error in the
observations, which is extremely small) 
\begin{align*}
&  \text{Likelihood in Problem \ref{lik3_prime}} \\ &= \P
                                                  \left\{\text{observed
                                                  data} \mid 
                                                  \theta, \sigma
                                                  \right\} \\
  &= \P\{X_1 \in [x_1 - \delta, x_1 + \delta], \dots, X_6 \in [x_6 -
    \delta, x_6 + \delta], X_1 \geq 25, \dots, X_5 \ge 25, X_6 < 25
    \mid \theta, \sigma \} 
\end{align*}
In other words, we are writing the probability that the first five
observations are all larger than 25 while the sixth observation is
smaller than 25, in addition to the exact values of the observations,
in the likelihood. But it is clear that these additional constraints
do not change the probability (as an example, just note that $\P(Z =
5) = \P(Z = 5, Z < 10)$. The additional restriction $Z < 10$ does not
affect the probability because it is already covered in $Z = 5$). Thus
\begin{align*}
&  \text{Likelihood in Problem \ref{lik3_prime}} \\ &= \P\{X_1 \in
                                                      [x_1 - \delta,
                                                      x_1 + \delta],
                                                      \dots, X_6 \in
                                                      [x_6 - 
    \delta, x_6 + \delta], X_1 \geq 25, \dots, X_5 \ge 25, X_6 < 25
                                                      \mid \theta, \sigma \} \\
  &= \P\{X_1 \in
                                                      [x_1 - \delta,
                                                      x_1 + \delta],
                                                      \dots, X_6 \in
                                                      [x_6 - 
    \delta, x_6 + \delta]    \mid \theta, \sigma \} \\
  &= \text{Likelihood in Problem \ref{lik3}}. 
\end{align*}
Since the likelihood is the same in both problems, Bayesian inference
for both problems will be the same (note the priors will be the same
as there is no reason to use different priors). Therefore, from the
Bayesian perspective, stopping rules  can be ignored for inferring $\theta$, 
because they do not affect the likelihood.

This example shows clearly that frequentist inference
violates the Likelihood Principle (the likelihood principle states
that ``all the evidence in a sample relevant to model parameters is
contained in the likelihood function''). See
\url{https://en.wikipedia.org/wiki/Likelihood_principle} for more
information on the likelihood principle.

On the other hand, Bayesian inference always satisfies the likelihood
principle (assuming that priors are the same), because data enters the
Bayesian posterior calculation only through the likelihood. 

Here is another example of violation of the likelihood principle in
frequentist inference. 

\section{Example 4: Coin Fairness Testing}

\begin{problem}
Suppose a coin is tossed 12 times leading to the outcome: TTTTHTHTTTTH
(this has 3 heads and 9 tails). What is your assessment of the
fairness of the coin?  
\end{problem}

\subsection{Frequentist Solution}

For the usual frequentist answer to this question, we assume that the
observed sequence of outcomes are the realization of random variables
$X_1, \dots, X_n$ (with $n = 12$) that are independently distributed
according to the $\text{Bin}(n, p)$ distribution for some unknown
$p$. We need to test the (null) hypothesis that $p = 0.5$ against,
say, the alternative $p < 0.5$. This can be 
done by calculating the $p$-value which is the probability (under the
assumption $p = 0.5$) of getting 3 or lower heads. The distribution of
the number of heads under the null distribution is $\text{Bin}(n,
0.5)$ so the $p$-value is
\begin{equation*}
 \left( \binom{12}{3} + \binom{12}{2} + \binom{12}{1}  + \binom{12}{0}
 \right) \frac{1}{2^{12}} = \frac{299}{4096} = 0.073 = 7.3\%
\end{equation*}
which does not lead to a rejection of the null hypothesis at the usual
5\% level.

In this $p$-value calculation, we
  implicitly assumed that the experiment 
consisted of tossing the coin 12 times where 12 was \textit{a priori}
chosen by the coin tosser. Consider now the alternative scenario where
the coin tosser wanted to toss the coin \textit{until the point where
  3 heads are 
observed}. Now for the same outcome, the $p$-value will change. Indeed 
now the random variable of interest will become $N = \text{number of tosses}$ and
the $p$-value will equal the probability of needing to toss the coin
12 or more times to get the 3 heads (assuming fairness). This is
calculated using the negative binomial distribution as: 
\begin{equation*}
 1 - \sum_{n = 3}^{11} \binom{n-1}{2} 2^{-n}  = \frac{134}{4096} =
 0.0327 = 3.27\%
\end{equation*}
and this leads to rejection of the null hypothesis at the $5\%$
level.


Note that the ``likelihood function'' is the same function $p^3(1 -
p)^9$ whether the sample size was predetermined or whether the coin
was tossed till 3 heads are observed. But the procedure obtained for
testing $p = 0.5$ has changed from the binomial to the negative
binomial case. This means that $p$-valued based frequentist inference
violates the Likelihood Principle. Here is a story from the
wikipedia article on the ``Likelihood Principle''  (see
\url{https://en.wikipedia.org/wiki/Likelihood_principle}) which puts
an interesting context to these numbers: 

\textit{Suppose a number of scientists are assessing the probability
of a certain outcome (which we shall call 'success') in experimental
trials. Conventional wisdom suggests that if there is no bias towards
success or failure then the success probability would be one
half. Adam, a scientist, conducted 12 trials and obtains 3 successes
and 9 failures. One of those successes was the 12th and last
observation. Then Adam left the lab.}
 
\textit{Bill, Adam's boss in the same lab, continued Adam's work and published
Adam's results, along with a significance test. He tested the null
hypothesis that $\theta$, the success probability, is equal to a half, versus
$\theta < 0.5$ . The probability that out of 12 trials, 3 or fewer
(i.e. more extreme) were successes, if $H_0$ is true, is $7.3\%$. Thus
the null hypothesis is not rejected at the 5\% significance level.}

\textit{Adam actually stopped immediately after 3 successes, because his boss Bill had
instructed him to do so. After the publication of the statistical
analysis by Bill, Adam realizes that he has missed a \textbf{later instruction}
from Bill to instead conduct 12 trials, and that Bill's paper is based
on this second instruction. Adam is very glad that he got his 3
successes after exactly 12 trials, and explains to his friend
Charlotte that by coincidence he executed the second
instruction. But Charlotte then explains to Adam that the $p$-value
should now be changed to $3.27\%$ and the result becomes significant
at the $5\%$ level. Adam is astonished to hear this.}

For more comments on the violation of the likelihood principle by
$p$-values, read \citet[Section 37.2]{MacKay}.

We shall look at the Bayesian solution to this problem in the next
lecture. 


\begin{thebibliography}{99}
    \bibitem[MacKay(2003)]{MacKay} MacKay, D. J. C., (2003)
  \textit{Information theory, inference, and learning
    algorithms}. Cambridge University Press, 2003.   

  \end{thebibliography}







\end{document}

%%% Local Variables:
%%% mode: latex
%%% TeX-master: t
%%% End:
